\section*{Limitations}
\label{sec:limitations}

Our human gold is pilot-scale: 5 papers and 55 cells with two annotators and a
single adjudication round, so the expert comparison is a preliminary signal rather
than a calibrated benchmark. The current scope is confined to
difference-in-differences identification; whether the rubric-plus-injection
methodology transfers to other quasi-experimental designs is untested. The rubric,
flaw taxonomy, and injector are co-designed, so synthetic detection rates are best
read as internal consistency under known threats rather than external
generalization; flaws authored by independent experts are future work. Relatedly,
flaw injection certifies the auditor against the \emph{injected} threat on a
single dimension only and inherits the taxonomy's blind spots, so we treat
detection and localization as necessary rather than sufficient, pairing them with
the false-alarm rate and the human-gold comparison. Evidence
grounding, not reasoning, is the dominant practical constraint: when retrieval returns
off-target passages the assessor cannot judge adequacy, which is why \sys{} emits an
explicit \texttt{unknown} state rather than a substantive verdict. Finally, \sys{}
is an auditor, not an oracle: it screens whether reported evidence supports a
study's identification assumptions and does not adjudicate the true causal effect,
which is never observed. Its outputs are meant to surface risk for expert review,
not to replace it.
